\documentclass{article}
\usepackage[UTF8]{ctex}
\usepackage{amsmath,amssymb,amsfonts}
\usepackage{enumitem}
\usepackage{geometry}
\geometry{a4paper, margin=1in}
% === 错题本排版优化（不参与题目标识） ===
\setlength{\parindent}{0pt}
\setlength{\parskip}{0.35em plus 0.10em minus 0.08em}
\setlength{\emergencystretch}{3em}
\setlist{itemsep=0.25em, topsep=0.40em, parsep=0pt}
\allowdisplaybreaks[2]
% === 错题本排版优化结束 ===
\begin{document}
\section*{河海大学2026年数学分析真题}
\section*{95．河海大学2026年数学分析试卷}

\subsection{题目1}
1．设 \(\lambda>1\) 及 \(\left\{a_{n}\right\}\) 是实数列，如果对所有 \(n, p \in \mathbb{N}^{+}\)都有 \(\left|a_{n+p}-a_{n}\right|<\frac{p}{n^{\lambda}}\) ，那么 \(\left\{a_{n}\right\}\) 是基本数列（Cauchy 列）。

\subsection{题目2}
2．设 \(x_{0}\) 是连续函数 \(f(x)\) 的极小值点，则存在 \(\delta>0\) ，使得 \(f(x)\) 在 \(\left(x_{0}-\delta, x_{0}\right]\) 上单调递减，在\(\left[x_{0}, x_{0}+\delta\right)\) 上单调递增．

\subsection{题目3}
3．设 \(f(x)\) 在 \([a, b]\) 上黎曼可积，则存在可导函数 \(F(x)\) ，使得 \(F^{\prime}(x)=f(x), x \in[a, b]\) ．

\subsection{题目4}
4．设实数列 \(\left\{a_{n}\right\}\) 满足 \(a_{n}>0(n=1,2, \cdots)\) ，且 \(\lim _{n \rightarrow \infty} a_{n}=0\) ，则级数 \(\sum_{n=1}^{\infty}(-1)^{n-1} a_{n}\) 收敛。

\subsection{题目5}
5．写出＂当 \(x \rightarrow+\infty\) 时，函数 \(f(x)\) 的极限为负无穷大＂的否定命题的分析表述．

\subsection{题目6}
6．叙述判断反常积分 \(\int_{a}^{+\infty} f(x) g(x) \mathrm{d} x\) 收敛的 Abel 判别法．

\subsection{题目7}
7．叙述判定函数项级数 \(\sum_{n=1}^{\infty} u_{n}(x)\) 关于 \(x\) 在集合 \(D\) 上一致收敛的 Cauchy 收敛原理．

\subsection{题目8}
8．叙述二元函数 \(z=z(x, y)\) 在点 \(\left(x_{0}, y_{0}\right)\) 处可微的定义．

\subsection{题目9}
9．计算极限 \(\lim _{x \rightarrow+\infty}\left(\sqrt[3]{x^{3}+2 x^{2}+1}-\sqrt[6]{x^{5}+3 x^{4}-2}\right)\) 。

\subsection{题目10}
10．计算不定积分 \(\int \frac{\mathrm{d} x}{\sqrt{\sin ^{7} x \cos x}}\) ．

\subsection{题目11}
11．求幂级数 \(\sum_{n=0}^{\infty} \frac{x^{n}}{4 n^{2}-1}\) 的收敛域及和函数。

\subsection{题目12}
12．设 \(\alpha>0\) ，且 \(f(x, y)\) 是 \(\mathbb{R}^{2}\) 上的可微函数．证明：对任意 \(t>0\) 及 \((x, y) \in \mathbb{R}^{2} \backslash\{(0,0)\}\) ，成立 \(f(t x, t y)=t^{\alpha} f(x, y)\) 的充要条件是对任意的 \((x, y) \in \mathbb{R}^{2}\) ，成立

\[
x \frac{\partial f}{\partial x}(x, y)+y \frac{\partial f}{\partial y}(x, y)=\alpha f(x, y) .
\]

\subsection{题目13}
13．计算曲面积分

\[
\iint_{\Sigma} \frac{x \mathrm{~d} y \mathrm{~d} z+y \mathrm{~d} z \mathrm{~d} x+z \mathrm{~d} x \mathrm{~d} y}{\left(x^{2}+y^{2}+z^{2}\right)^{3}} .
\]

其中 \(\sum\) 是曲面 \(9-z=\sqrt{(x-1)^{2}+(y-2)^{2}}(z \geq 0)\) ，方向取上侧。

\subsection{题目14}
14．设 \(f(x)\) 和 \(g(x)\) 在 \([a, b]\) 上二阶可导，且 \(g^{\prime \prime}(x) \neq 0, f(a)=f(b)=g(a)=g(b)=0\) ．

\begin{itemize}
  \item[（1）] 证明：对任意的 \(x \in(a, b)\) ，都有 \(g(x) \neq 0\) ．
  \item[（2）] 证明：存在 \(\xi \in(a, b)\) ，使得 \(\frac{f(\xi)}{g(\xi)}=\frac{f^{\prime \prime}(\xi)}{g^{\prime \prime}(\xi)}\) ．
\end{itemize}

\subsection{题目15}
15．设 \(\alpha>0\) ．

\begin{itemize}
  \item[（1）] 证明：反常积分 \(\int_{1}^{+\infty} \frac{\mathrm{d} x}{\left(1+x^{2} y^{2}\right) x^{\alpha}}\) 关于 \(y\) 在 \((0,+\infty)\) 上内闭一致收敛．
  \item[（2）] 证明：\(I(y)=\int_{1}^{+\infty} \frac{\arctan (x y)}{x^{\alpha+1}} \mathrm{~d} x\) 是 \((0,+\infty)\) 上的可导函数，且满足方程
\[
I^{\prime}(y)-\alpha I(y)+\arctan y=0 .
\]
\end{itemize}

\subsection{题目16}
16．设 \(f(x)\) 在 \([a, b]\) 上连续，且满足 \(0<m \leq f(x) \leq M\) ，证明：

\[
1 \leq \frac{1}{(b-a)^{2}} \int_{a}^{b} f(x) \mathrm{d} x \int_{a}^{b} \frac{\mathrm{~d} x}{f(x)} \leq \frac{(m+M)^{2}}{4 m M}
\]

\subsection{题目17}
17．设 \(a_{n} \neq 0(n=1,2, \cdots)\) ，对级数 \(\sum_{n=1}^{\infty} a_{n}\) 加括号得到

\[
\left(a_{1}+a_{2}+\cdots+a_{n_{1}}\right)+\left(a_{n_{1}+1}+a_{n_{1}+2}+\cdots+a_{n_{2}}\right)+\cdots+\left(a_{n_{k-1}+1}+a_{n_{k-1}+2}+\cdots+a_{n_{k}}\right)+\cdots
\]

其中每个括号内的 \(a_{n}\) 具有相同的符号，相邻括号内 \(a_{n}\) 的符号不同。记 \(n_{0}=0\) ，且

\[
b_{k}=\sum_{j=n_{k-1}+1}^{n_{k}} a_{j}(k=1,2, \cdots) .
\]

（1）证明：\(\sum_{n=1}^{\infty} a_{n}\) 和 \(\sum_{n=1}^{\infty} b_{n}\) 具有相同的敛散性．\\
（2）证明：\(\sum_{n=1}^{\infty} \frac{(-1)^{[\sqrt{n}]}}{n}\) 是收敛的，其中 \([x]\) 是取整函数，表示不超过 \(x\) 的最大整数。
\end{document}
